\documentclass[10pt,a4paper]{amsart}
\usepackage[margin=19mm]{geometry}
\usepackage{amsmath,amssymb,mathtools}
\usepackage[scr=rsfs,cal=euler]{mathalfa}
\usepackage{xfrac}
\usepackage[unicode,hidelinks,bookmarks=false]{hyperref}
\newcommand{\ie}{i.e.\ }
\newcommand{\cf}{cf.\ }
\newcommand{\resp}{respectively\ }
\newcommand{\Subsection}{\subsection}
\providecommand{\nobreakdash}{\nobreak\hbox{-}}
\setlength{\emergencystretch}{2em}
\multlinegap0pt
\usepackage{amssymb}
\usepackage{mathtools}
\usepackage{enumerate}
\usepackage{tikz}
\usepackage{xcolor}
\newtheorem{lemma}{Lemma}
\newcommand{\Z}{\mathbb{Z}}
\title{Stably exotic fillings of 3-manifolds}
\author{Daniel Kasprowski, Patrick Orson, Mark Powell, Arunima Ray}
\date{Source: arXiv:2608.23523v1 (CC BY 4.0)}
\begin{document}
\maketitle

\noindent\emph{Source and licence.} Stably exotic fillings of 3-manifolds. Version arXiv:2608.23523v1 (CC BY 4.0), distributed by arXiv under the Creative Commons Attribution 4.0 International licence, http://creativecommons.org/licenses/by/4.0/, as recorded in the arXiv licence metadata for this exact version and shown on the abstract page https://arxiv.org/abs/2608.23523v1. Full licence notice: \texttt{licenses/arxiv-2608.23523-CC-BY-4.0.txt}.\par\smallskip

\noindent\textit{Editorial scope.} Counted unit: the article's lemma lemma:lagrangian-complement (source0.tex lines 916--918). The article's complete original proof of it is delivered with it (source0.tex lines 920--936, one \texttt{proof} environment of 17 lines). No mathematics is written, repaired or re-derived: the statement and the proof are the article's own lines, in the article's own notation, and no retained line is substituted or re-flowed.  No further source-local context is needed: every cross-reference of every retained line resolves inside the two delivered spans.  Nothing was omitted from any retained span, so the package carries no omission record.  No retained line contains a citation, so the wrapper carries no bibliography and no bibliography entry.  No retained line contains a figure environment, an image-inclusion command or drawing code, so no image is delivered and the package declares no asset dependency.  Editorial formatting only, in addition: an AMS \texttt{amsart} preamble that restates the article's own declarations and macros used by the retained lines, namely the environments \texttt{lemma} on separate counters, together with the standard packages those lines need.  The retained environments print in this document as Lemma 1 in order of appearance, whereas the article prints the same items with the numbering of its own full document.  The complete native source of the article as distributed in the arXiv e-print for this exact version is bundled unchanged as \texttt{sources/208415/source0.tex}.
\begin{lemma}\label{lemma:lagrangian-complement}
    Let $(P,\psi)$ be a hyperbolic form over $\Z/2$ of dimension $2g$ and let $V$ be a subspace of $P$ of dimension at least $g$. Then there is a Lagrangian $K \subseteq P$ such that $K +V = P$. 
\end{lemma}

\begin{proof}
By passing to a $g$-dimensional subspace of $V$, without loss of generality we can assume that~$V$ is rank $g$. Writing $j\colon V\to P$ for the inclusion, we have that $V^\perp$ is the kernel of the surjective map $j^*\circ \mathrm{ad}(\psi)\colon P\to V^*$, so  $\dim V^\perp=\dim P-\dim V^*=g = \dim V$. Let $Z:=V\cap V^\perp$ and choose direct sum decompositions $V=U\oplus Z$ and $V^\perp=U'\oplus Z$, so that $V+V^\perp=Z\oplus U\oplus U'$. Noting that $Z$ is the radical for the restriction of $\psi$ to both $V$ and to $V^\perp$, we have that $\psi$ restricts to a nonsingular form on both $U$ and $U'$. As $U$ and $U'$ have the same dimension, we may thus choose an isomorphism of nonsingular alternating forms $f\colon (U,\psi|_U)\to (U\,\psi|_{U'})$. Such an isomorphism determines a Lagrangian $\Delta_f:=\{(y,f(y))\,|\,y\in U\}\subseteq U\oplus U'$.  By the same argument as the second sentence of the proof, $\dim(U \oplus U')^{\perp} = \dim P - \dim(U \oplus U')$. 
The restriction of $\psi$ to $U\oplus U'$ is nonsingular, so $(U\oplus U')\cap (U\oplus U')^\perp=\{0\}$. Combining the two previous sentences, we have 
\begin{equation}\label{eqn:decomposition-of-P}
P=(U\oplus U')\oplus (U\oplus U')^\perp.
\end{equation}
We will now extend $\Delta_f$ to a full Lagrangian of $(P,\psi)$. Note that $Z\subseteq (U\oplus U')^\perp$. 
Since $\dim U + \dim Z = \dim V = g = \dim V^{\perp} = \dim U' + \dim Z$ we have that $2 \dim Z = 2g-\dim U - \dim U'$. Also $2g = \dim P = \dim U + \dim U' + \dim (U\oplus U')^{\perp}$, and hence $\dim (U\oplus U')^{\perp} = 2 \dim Z$, i.e.\ $Z\subseteq (U\oplus U')^\perp$ is half-rank.  Moreover we already noted that $\psi$ vanishes identically on $Z$, so $Z$ is isotropic.  
Since the decomposition in \eqref{eqn:decomposition-of-P} is by definition orthogonal with respect to $\psi$, and $\psi$ is nonsingular, $\psi$ restricts to a nonsingular form on $(U\oplus U')^\perp$.  It follows that $Z$ is a Lagrangian for the restriction in $(U\oplus U')^\perp$. Moreover $Z$ admits a complementary Lagrangian $L$,  so that $Z\oplus L= (U\oplus U')^\perp$.  
To find one,  start with a complementary summand $L'$. Since $\psi|$ is nonsingular, $Z \to (L')^*$ is an isomorphism, and hence by adding elements of $Z$ to the elements of a basis for~$L'$, we can change $L'$ to a Lagrangian $L$ that is still complementary to $Z$. 

For any complementary Lagrangian $L$ to $Z$ in $(U\oplus U')^\perp$, the submodule
\[
K:=\Delta_f\oplus L\subseteq (U\oplus U')\oplus (U\oplus U')^\perp
\]
is a Lagrangian, and is such that $K+V=P$.
\end{proof}

\end{document}
